% from ega0-20-fr.tex
\oldpage[0\textsubscript{IV-1}]{126}
(i) Il est immédiat que $\varphi$ (avec les notations de
\sref{0.20.4.6-fr}) est nécessairement l'application
$z\mapsto(\varepsilon(z),z-p_1(\varepsilon(z)))$, et l'isomorphisme
réciproque l'application $(b,x)\mapsto p_1(b)+x$ ; ces deux applications
étant continues, cela prouve la première assertion. On notera que cela
entraîne que la topologie de $B$ s'identifie (par $p_1$) au \emph{quotient}
de la topologie de $B\otimes_A B$ par $\mathfrak{J}_{B/A}$.

(ii) Le fait que $d_{B/A}$ soit une $A$-dérivation de $B$ résulte de la
définition \sref{0.20.4.6-fr} et de \sref{0.20.1.1-fr}.

Pour prouver la propriété universelle, rappelons que
$\operatorname{D\acute{e}r.\ cont}_A(B,L)$ s'identifie canoniquement à
l'ensemble $G$ des homomorphismes continus de $A$-algèbres
$u:B\to D_B(L)$ tels que le composé
$B\xrightarrow{u}D_B(L)\xrightarrow{q}B$ soit l'identité (où $q(b,0)=b$ est
la projection ; cf. \sref{0.20.1.6-fr} et \sref{0.20.3.2-fr}). D'autre part,
grâce à l'isomorphisme $\varphi$,
$\operatorname{Hom.\ cont}_B(\Omega^1_{B/A},L)$ s'identifie canoniquement à
l'ensemble des homomorphismes continus de $B$-algèbres
$v:\mathrm{P}^1_{B/A}\to D_B(L)$ tels que le composé
$\mathrm{P}^1_{B/A}\xrightarrow{v}D_B(L)\xrightarrow{q}B$ soit
l'augmentation $\varepsilon$. Comme on a $p_2=p_1-d_{B/A}$ par définition,
tout revient à prouver que tout $u\in G$ se factorise en
\[
\xymatrix{
  B \ar[r]^u \ar[d]_{p_2} & D_B(L) \\
  \mathrm{P}^1_{B/A} \ar[ur]_v &
}
\]
où $v$ est un $B$-homomorphisme continu. Or, on a déjà un homomorphisme
continu de $A$-algèbres $j:b\mapsto(b,0)$ de $B$ dans $D_B(L)$, qui
appartient à $G$ ; par définition du produit tensoriel topologique
d'algèbres topologiques ($\mathbf{0_I}$, 7.7.6), il existe donc un
$A$-homomorphisme continu d'algèbres
$w:B\otimes_A B\to D_B(L)$ rendant commutatif le diagramme
\[
\xymatrix{
  B\otimes_A B \ar[dr]^w & B \ar[l]_{j_1} \ar[d]^j \\
  B \ar[u]^{j_2} \ar[r]_u & D_B(L)
}
\]
On a donc par définition
$w(b\otimes1-1\otimes b)=j(b)-u(b)\in L$, et en vertu de
\sref{0.20.4.4-fr}, cela entraîne $w(\mathfrak{J})\subset L$ donc
$w(\mathfrak{J}^2)=0$ ; par suite $w$ se factorise en
\[
  B\otimes_A B\to\mathrm{P}^1_{B/A}\xrightarrow{v}D_B(L)
\]
où $v$ est un homomorphisme continu de $A$-algèbres ; d'ailleurs, comme
$v\circ p_1=j$ est un homomorphisme de $B$-algèbres, il en est de même de
$v$ par définition de la structure de $B$-algèbre de
$\mathrm{P}^1_{B/A}$ ; comme on a par définition $u=v\circ p_2$, cela
termine la démonstration.
\end{proof}

\begin{theorem}[20.4.9]
\label{0.20.4.9-fr}
Supposons que $B$ soit une $A$-algèbre topologique formellement lisse. Alors
le $B$-module topologique $\Omega^1_{B/A}$ est \emph{formellement projectif}.
\end{theorem}

\begin{proof}
En effet, l'hypothèse entraîne que $B\otimes_A B$ (muni de la structure de
$B$-algèbre topologique définie par $j_1$) est une $B$-algèbre topologique
formellement lisse \sref{0.19.3.5-fr}, (iii), et par suite aussi une
$A$-algèbre topologique formellement lisse \sref{0.19.3.5-fr}, (ii) ; comme
$B$ est topologiquement isomorphe à la $A$-algèbre quotient
$(B\otimes_A B)/\mathfrak{J}$ et est une $A$-algèbre formellement lisse, la
conclusion résulte de \sref{0.19.5.3-fr}.
\end{proof}

